\section{Web Security Attacks}
\label{sec::chp2::security-attacks}

%TODO: OSI layer attacks
%\begin{table}[h!]
\footnotesize
\centering
\begin{tabular}{|c|l|}
\hline
\textbf{OSI Layer} & \textbf{Attack Types} \\
\hline
\multirow{5}{*}{Physical (Layer 1)} & Wiretapping \\
 & Eavesdropping \\
 & Cable cutting \\
 & Electromagnetic interference (EMI) \\
 & Power fluctuations \\
\hline
\multirow{5}{*}{Data Link (Layer 2)} & MAC spoofing \\
 & ARP spoofing \\
 & VLAN hopping \\
 & STP (Spanning Tree Protocol) attacks \\
 & CDP (Cisco Discovery Protocol) attacks \\
\hline
\multirow{7}{*}{Network (Layer 3)} & IP spoofing \\
 & ICMP flooding \\
 & Smurf attack \\
 & Ping of Death \\
 & Teardrop attack \\
 & IP fragmentation attacks \\
 & Routing attacks (e.g., BGP hijacking) \\
\hline
\multirow{5}{*}{Transport (Layer 4)} & SYN flooding \\
 & TCP session hijacking \\
 & UDP flooding \\
 & Port scanning \\
 & Firewalk attack \\
\hline
\multirow{4}{*}{Session (Layer 5)} & Session hijacking \\
 & Session fixation \\
 & NetBIOS attacks \\
 & RPC (Remote Procedure Call) attacks \\
\hline
\multirow{3}{*}{Presentation (Layer 6)} & ASN.1 (Abstract Syntax Notation One) attacks \\
 & Data encoding attacks \\
 & Format string attacks \\
\hline
\multirow{12}{*}{Application (Layer 7)} & Cross-site scripting (XSS) \\
 & SQL injection \\
 & Remote code execution \\
 & Buffer overflow \\
 & Denial of Service (DoS) \\
 & Man-in-the-Middle (MitM) attacks \\
 & Phishing attacks \\
 & Malware (viruses, worms, Trojans) \\
 & Brute-force attacks \\
 & Directory traversal attacks \\
 & Injection flaws \\
 & Insecure direct object references \\
 & Broken authentication and session management \\
\hline
\end{tabular}
\caption{Different types of attacks against each of the OSI layers}
\label{tab:osi-layer-attacks}
\end{table}

This section provides an overview of malicious techniques used for unauthorised exploitation of web applications (not exclusively) and exfiltration of user secrets. We will refer to some of these attacks through the thesis. The main purpose is for this section to act as a reference and guide. This is a comprehensive version of the OWASP's list of attacks. https://owasp.org/www-community/attacks/

Seldom is a single threat vector used exclusively in the real world. Often, multiples areas of weaknesses are exploited to complete a mission. Therefore, having a broad perspective of these attacks and their relationships with each other is essential.


\subsection{Client Side Attacks}
\label{chp2::client-side-attacks}

\textbf{\textsc{Credential Stuffing Attack.}} Attackers use Credential Stuffing Attack (CSA) to gain unauthorised account and system access. They do this by using compromised login information from previous data breaches. CSA owes its success to the inherent weaknesses in the Password. Such weaknesses include: re-using passwords across different systems and using weak passwords.

The increasing number of data breaches has made this form of attack viable. The heap of data that is often dumped on the internet increases the feasibility of CSA. This accelerating threat has led to the likes of SEC issuing warnings to the industry \cite{noauthor_secgov_2020}.

CSA is one of the many consequences of dumped private data. Others include identity fraud and attempted blackmailing by emailing the victims of the data breaches \cite{noauthor_scam_2020}.

\textbf{\textsc{Man-in-the-Browser.}} MitB is usually a Trojan that infects a system and its applications (typically a browser). Once it infects a browser or an operating-system (OS) it can read network traffic, inject code into web pages, change HTTP headers, and generally manipulate the web page for the benefit of the attacker \cite{noauthor_man---browser_2020}. MitB Trojans are also referred to as \emph{Financial Trojans} because (typically) their aim is commit fraud and steal money from their victims. Such malicious actions are done by form grabbing banking details, personal details and changing transaction details (if the user is using online banking to transfer money).

One of the oldest and most famous MitB Trojan is ZeuS \cite{noauthor_zeus_2010}, also known as Zbot. First discovered in 2006, it has caused vast amount of financial damages. Its creator reportedly retired in 2010 and released \cite{noauthor_zeus_2011} its source code to ZeuS' competitor, SpyEye \cite{noauthor_trojan.spyeye_nodate}.

Despite improvements in browser and OS security, criminals are finding
creative ways in penetrating user systems in order to commit malicious
activities.

Mobile phones are not safe either. A variant of ZeuS known as Zitmo \cite{etaher_zeus_2015} or The-Zeus-in-the-Mobile targets mobile phone
devices.

Form grabbing is a type of MitB that has one focus: to skim and steal
contents of forms that the user is filling out. Form grabbers are often
malicious JavaScript code running on compromised websites.

MitB malware is prevalent in the financial sector and payment gateway
systems. Using this method, the attackers maximise their financial gain.
A notable example of MitB---in the form of form grabbing--- is the
British Airways' incident: the attackers were able to manipulate a
third-party JavaScript running on British Airways' payment page. This
manipulation allowed them to scrape payment data as they were entered on
the site by British Airways' customers \cite{klijnsma_british_2018}.

\textbf{\textsc{Password Spraying Attack.}}

%   \begin{figure}[H]
%           \centering
%           \textbf{Authentication with explicit link to deeID}
%           \begin{sequencediagram}
%               \centering
%               \newinst[1]{A}{User}
%               \newinst[1]{B}{Client}
%               \newinst[1]{C}{Server}
%               \newinst[1]{D}{Attacker}
%               \mess{D}{Infect JS Skimmer}{C}
%               \mess{B}{GET /payment}{C}
%               \mess{C}{Respond /payment}{B}
%               \mess{A}{Type payment data $= d$}{B}
%               \mess{A}{Press Submit}{B}
%               \mess{B}{POST /m, body $= d$}{D}
%               \mess{B}{POST /payment, body $= d$}{C}
%           \end{sequencediagram}
%           \label{seq:basicmitb}
%           \caption{Sequence of a basic MitB Skimmer.}
%       \end{figure}

\begin{enumerate}
\def\labelenumi{\arabic{enumi}.}
\item
  \textbf{Clickjacking or UI Redress Attack}: This is a deceptive technique where a
  malicious website or attacker tricks a user into clicking on something
  different from what they perceive. It is typically done by overlaying
  transparent elements on a legitimate webpage, making users perform
  unintended actions without their knowledge.
\item
  \textbf{Cross Frame Scripting (XFS)}:  This is an attack that combines malicious JavaScript with an iframe that loads a legitimate page in an effort to steal data from an unsuspecting user. This attack is usually only successful when combined with social engineering. An example would consist of an attacker convincing the user to navigate to a web page the attacker controls. The attacker’s page then loads malicious JavaScript and an HTML iframe pointing to a legitimate site. Once the user enters credentials into the legitimate site within the iframe, the malicious JavaScript steals the keystrokes.
\item
  \textbf{Cross Site History Manipulation (XSHM)}: XSHM is an attack
  that manipulates the browser's history to change the URLs visited by a
  user. This can lead to confusing and potentially harmful navigation,
  exploiting the same-origin policy of web browsers.
\item
  \textbf{Cross Site Tracing}: Cross Site Tracing is an attack that
  takes advantage of HTTP TRACE requests to steal sensitive information,
  often through cross-site scripting or other vulnerabilities. It can
  expose cookies or credentials to attackers.
\item
  \textbf{Mobile code invoking untrusted mobile code}: This refers to
  situations where mobile code (code executed on a mobile device, like a
  smartphone) invokes other untrusted code, which could potentially lead
  to security risks. It's a concern in the context of mobile app
  development.
\item
  \textbf{Mobile code non-final public field}: This relates to the
  exposure of non-final public fields in mobile code, which can be
  manipulated or tampered with by external parties, potentially leading
  to security vulnerabilities.
\item
  \textbf{Mobile code object hijack}: In this attack, malicious mobile
  code can hijack or take control of objects or resources in a mobile
  environment, potentially leading to unauthorized access or
  manipulation.
\item
  \textbf{Qrljacking}: Qrljacking is a type of attack where an attacker
  tricks users into scanning a QR code that leads to a malicious website
  or action, rather than the intended destination.
\item
  \textbf{Reflected DOM Injection}: Reflected DOM Injection is a type of
  cross-site scripting (XSS) attack where the injected script
  manipulates the Document Object Model (DOM) of a web page. This can
  lead to various security issues and data theft.
\item
  \textbf{Reverse Tabnabbing}: Reverse Tabnabbing is a variant of
  phishing that occurs when a malicious website changes the content of a
  previously opened tab or window, making it appear legitimate while
  tricking the user into revealing sensitive information.
\item
  \textbf{Cross-Site Scripting (XSS)}: XSS is a widespread attack where
  an attacker injects malicious scripts into web content that is then
  executed by unsuspecting users. It can result in data theft, session
  hijacking, or defacement of websites. It can occur in various
  contexts, including in subtitles, if not properly sanitized.
\item
  \textbf{Cross Site Request Forgery (CSRF)}: CSRF is an attack where an
  attacker tricks a user into unknowingly making an unwanted request to
  a different website on which the user is authenticated, potentially
  causing unauthorized actions.
\item
  \textbf{Drive-By Downloads}: Drive-By Downloads are attacks that
  automatically download and install malicious software onto a user's
  device when they visit a compromised or malicious website, often
  without their consent.
\item
  \textbf{Browser Fingerprinting}: Browser fingerprinting is a technique
  used to collect unique information about a user's web browser and
  device. This information can be used to track users without their
  knowledge, compromising privacy.
\item
  \textbf{Content Spoofing - Form Spoofing}: Content spoofing,
  particularly form spoofing, involves manipulating the content of web
  forms to deceive users into providing sensitive information to a
  malicious website or attacker.
\item
  \textbf{Malvertising}: Malvertising involves delivering malicious
  software or content through online advertisements. Users may
  unknowingly encounter malware by clicking on these ads.
\item
  \textbf{Session Fixation}: Session fixation is an attack where an
  attacker sets or fixes a user's session identifier, potentially
  allowing them to impersonate the victim after login.
\item
  \textbf{Malicious Browser Extensions}: Malicious browser extensions
  are browser add-ons or plugins that appear harmless but can steal
  data, inject ads, or compromise user privacy.
\end{enumerate}

%TODO: Server side attacks
%\subsection{Server-Side Attacks}

\begin{itemize}
\item \textbf{Binary Planting:} Uploading and executing unauthorized binary files on the server, usually a specific location. Binary planting can be used to escalate privileges. An example of this is the Stuxnet virus using the  Microsoft Windows Print Spooler Service Remote Code Execution Vulnerability (CVE-2010-2729 / BID 43073). Also known as \textbf{DLL Planting} or \textbf{DLL Hijacking}.
\item \textbf{SQL Injection}
\begin{itemize}
    \item Blind SQL Injection: Extracting data from a database by sending queries and observing the server's response time or behavior.
\end{itemize}
\item Blind XPath Injection: Extracting data from an XML document by sending queries and observing the server's response time or behavior.
\item Buffer Overflow via Environment Variables: Manipulating environment variables to cause a buffer overflow and execute arbitrary code.
\item Buffer Overflow Attack: Overwriting memory beyond the bounds of a buffer to corrupt data or execute arbitrary code.
\item \textbf{CORS}
\begin{itemize}
    \item CORS Origin Header Scrutiny: Exploiting misconfigured Cross-Origin Resource Sharing (CORS) headers to access restricted resources.
    \item CORS Request Preflight Scrutiny: Exploiting misconfigured CORS preflight requests to bypass security checks.
\end{itemize}
\item \textbf{File Injection:}
\begin{itemize}
    \item CSS 
    \item CSV Injection: Injecting malicious payloads into CSV files to execute arbitrary code or perform unauthorized actions.
\end{itemize}

\item Cache Poisoning: Manipulating cached data to serve malicious content to users.
\item Cash Overflow: Manipulating financial calculations to cause integer overflows and perform unauthorized transactions.
\item Code Injection: Injecting malicious code into a vulnerable application to execute arbitrary commands or manipulate data.
\item Command Injection: Executing arbitrary commands on the server by injecting them into vulnerable input fields.
\item Comment Injection Attack: Injecting malicious comments or code into an application's data store or output.
\item Content Spoofing: Tricking users into believing that malicious content is legitimate by mimicking the appearance of a trusted website.
\item Custom Special Character Injection: Injecting special characters into input fields to bypass validation or cause unintended behavior.
\item Denial of Service (DoS): Overwhelming a server with requests to exhaust its resources and make it unavailable to users.
\item Direct Dynamic Code Evaluation - Eval Injection: Injecting malicious code into an application's eval() function to execute arbitrary commands.
\item Embedding Null Code: Injecting null bytes into input fields to bypass validation or manipulate file paths.
\item Execution After Redirect (EAR): Executing malicious code after a redirect to bypass security checks.
\item Forced Browsing: Accessing restricted pages or resources by guessing or manipulating URLs.
\item Form Action Hijacking: Modifying the action attribute of HTML forms to submit data to an attacker-controlled server.
\item Format String Attack: Exploiting format string vulnerabilities to read or write arbitrary memory locations.
\item Full Path Disclosure: Revealing the full path of files or directories on the server.
\item Function Injection: Injecting malicious code into a vulnerable function to execute arbitrary commands or manipulate data.
\item HTTP Response Splitting: Injecting newline characters into HTTP headers to split the response and insert malicious content.
\item LDAP Injection: Injecting malicious LDAP queries to extract sensitive information or manipulate directory services.
\item Log Injection: Injecting malicious payloads into application logs to execute arbitrary code or manipulate log data.
\item Server-Side Includes (SSI) Injection: Injecting SSI directives to execute arbitrary commands or include malicious files.
\item Server-Side Request Forgery (SSRF): Tricking a server into making unauthorized requests to internal or external resources.
\item Session Prediction: Guessing or predicting valid session identifiers to hijack user sessions.
\item Session Fixation: Forcing a user to use a predefined session identifier to hijack their session.
\item Session Hijacking Attack: Stealing or predicting session identifiers to take over user sessions.
\item Setting Manipulation: Modifying application settings or configuration files to alter its behavior or gain unauthorized access.
\item Special Element Injection: Injecting special elements or tags into input fields to bypass validation or manipulate the application's behavior.
\item Unicode Encoding: Exploiting Unicode encoding vulnerabilities to bypass input validation or inject malicious payloads.
\item Web Parameter Tampering: Manipulating URL parameters or form data to perform unauthorized actions or bypass security checks.
\item Windows ::DATA Alternate Data Stream: Hiding malicious files or data in alternate data streams on Windows systems.
\item XML External Entity (XXE) Injection: Exploiting XML parsers to access unauthorized files or network resources.
\item XPath Injection: Injecting malicious XPath queries to extract sensitive information from XML documents.
\item Cross-Site Request Forgery (CSRF): Tricking users into performing unintended actions on a web application in which they are authenticated.
\item Web Service Amplification Attack: Exploiting web services to amplify denial of service attacks by sending large amounts of data.
\end{itemize}


\hypertarget{other-attacks}{%
\subsubsection{Other Attacks:}\label{other-attacks}}

\begin{itemize}
\item
  Brute Force Attack
\item
  Cross-User Defacement
\item
  Forced browsing
\item
  Form action hijacking by Robert Gilbert (amroot)
\item
  Man-in-the-browser attack
\item
  Manipulator-in-the-middle attack
\item
  Parameter Delimiter
\item
  Path Traversal
\item
  Regular expression Denial of Service - ReDoS by Adar Weidman
\item
  Repudiation Attack
\item
  Resource Injection
\item
\begin{verbatim}
* **Inband**. The same channel is used to inject SQL code and extract sensitive data.
\end{verbatim}
\item
\begin{verbatim}
* **Out-of-band.**  Data is retrieved using a different channel.
\end{verbatim}
\item
\begin{verbatim}
* **Inferential.**  Gaining information from behaviour server based on sent requests.
\end{verbatim}
\item
\begin{verbatim}
* Tools:
\end{verbatim}
\item
\begin{verbatim}
    * mieliekoek.pl
\end{verbatim}
\item
\begin{verbatim}
    * wpoison
\end{verbatim}
\item
\begin{verbatim}
    * sqlmap
\end{verbatim}
\item
\begin{verbatim}
    * wapiti
\end{verbatim}
\item
\begin{verbatim}
    * w3af
\end{verbatim}
\item
\begin{verbatim}
    * paros
\end{verbatim}
\item
\begin{verbatim}
    * sqid
\end{verbatim}
\item
\begin{verbatim}
* User input
\end{verbatim}
\item
\begin{verbatim}
* Cookies
\end{verbatim}
\item
\begin{verbatim}
* HTTP Headers
\end{verbatim}
\item
\begin{verbatim}
* Second-order SQL Injections
\end{verbatim}
\item
  Traffic flood
\item
  Trojan Horse
\item
  Unicode Encoding
\item
  Windows ::DATA Alternate Data Stream
\end{itemize}